\documentclass{article}

\textwidth=125mm
\textheight=200mm
\linespread{1.04}


\title{\textsc{Optimistic BCO}}
\usepackage{marginnote}
\reversemarginpar % switch default margin to the left
\usepackage{framed}
\usepackage{xcolor}
\usepackage{nicefrac}
\usepackage{pgfplots}
\usepackage{floatrow}
\usepackage{mathtools}
\usepackage{amsthm}
\usepackage{listings}
\usepackage{tikz}
\usepackage{amssymb}
\usepackage[plain]{algorithm}
\lstset{emph={try,catch,throw,def,or,to,for,then,end,for,if,args,else,return},emphstyle=\color{blue!100!black}\textbf}
\lstset{mathescape=true,escapechar=\&,numbers=left,xleftmargin=3mm,linewidth=0.47\textwidth,framextopmargin=0pt,framexbottommargin=0pt,aboveskip=-1pt,belowskip=-1pt}
\usepackage{mdframed}
\usepackage{bm}
\usepackage[bf]{caption}
\usepackage{graphicx}
\usepackage{natbib}
\usepackage{xspace}
\usepackage[hypertexnames=false]{hyperref} % add backlinks to the end of bibliography items
\usepackage[capitalise]{cleveref}
\usepackage{autonum}

\definecolor{dkblue}{cmyk}{1,.54,.04,.19} 
\theoremstyle{plain}
\newtheorem{theorem}{Theorem}
\newtheorem{optprob}[theorem]{Optimisation Problem}
\newtheorem{lemma}[theorem]{Lemma}
\newtheorem{proposition}[theorem]{Proposition}
\newtheorem{corollary}[theorem]{Corollary}
\newtheorem{claim}[theorem]{Claim}
\theoremstyle{definition}
\newtheorem{definition}[theorem]{Definition}
\newtheorem{question}[theorem]{Question}
\newtheorem{fact}[theorem]{Fact}
\newtheorem{assumption}[theorem]{Assumption}
\newtheorem{conjecture}[theorem]{Conjecture}
\theoremstyle{plain}
\newtheorem{exampleinner}[theorem]{Example}
\newtheorem{remark}[theorem]{Remark}

% Boxed example environment with light gray background and shadow
\mdfdefinestyle{examplestyle}{%
    backgroundcolor=gray!10,
    linewidth=0.5pt,
    linecolor=gray!50,
    shadow=true,
    shadowsize=3pt,
    shadowcolor=gray!30,
    innertopmargin=8pt,
    innerbottommargin=8pt,
    innerleftmargin=10pt,
    innerrightmargin=10pt,
    skipabove=10pt,
    skipbelow=10pt
}
\newenvironment{example}[1][]{%
    \begin{mdframed}[style=examplestyle]%
    \begin{exampleinner}[#1]%
}{%
    \end{exampleinner}%
    \end{mdframed}%
}



\newcommand{\ceil}[1]{\left\lceil #1 \right\rceil}
\newcommand{\R}{\mathbb R}
\newcommand{\Rd}{\mathbb{R}^d}
\newcommand{\argmin}{\operatornamewithlimits{arg\,min}}
\newcommand{\argmax}{\operatornamewithlimits{arg\,max}}
\newcommand{\ip}[1]{\langle #1 \rangle}
\newcommand{\Reg}{\mathfrak{R}}
\newcommand{\BReg}{\mathfrak{BR}}
\newcommand{\norm}[1]{\Vert #1 \Vert}
\newcommand{\E}{\mathbb E}
\newcommand{\Z}{\mathbb Z}
\newcommand{\K}{\mathcal K}
\newcommand{\cU}{\mathcal U}
\newcommand{\cN}{\mathcal N}
\newcommand{\cT}{\mathcal T}
\newcommand{\PP}{\mathbb P}
\newcommand{\Var}{\mathbb V}
\newcommand{\cE}{\mathcal E}
\newcommand{\cK}{\mathcal K}
\newcommand{\cH}{\mathcal H}
\newcommand{\cF}{\mathcal F}
\newcommand{\cO}{\mathcal O}
\newcommand{\cS}{\mathcal S}
\newcommand{\cC}{\mathcal C}
\newcommand{\cL}{\mathcal L}
\newcommand{\cG}{\mathcal G}
\newcommand{\cA}{\mathcal A}
\newcommand{\cD}{\mathcal D}
\newcommand{\cP}{\mathcal P}
\newcommand{\BS}{\mathbb S}
\newcommand{\zeros}{ \bm 0}
\newcommand{\bbP}{\mathbb P}
\newcommand{\polylog}{\operatorname{polylog}}
\newcommand{\poly}{\operatorname{poly}}
\newcommand{\const}{\operatorname{const}}
\newcommand{\interior}{\operatorname{int}}
\newcommand{\conv}{\operatorname{conv}}
\newcommand{\diam}{\operatorname{diam}}
\renewcommand{\d}[1]{\operatorname{d}\!#1}
\newcommand{\abovelabel}[1]{\stackrel{\raisebox{0.5mm}{\textrm{\tiny \color{red} #1}}}}
\newcommand{\ta}{\tilde{a}}
\newcommand{\tP}{P'}
\newcommand{\tA}{A'}
\newcommand{\fmd}{\texttt{MD}}
\newcommand{\tfmd}{\widetilde{\texttt{MD}}}

\newcommand{\mdalg}{\textsf{MD}\xspace}
\newcommand{\smdalg}{\textsf{SMD}\xspace}
\newcommand{\expbyopt}{\textsf{ExpByOpt}\xspace}
\newcommand{\olo}{\texttt{OLO}\xspace}
\newcommand{\blo}{\texttt{BLO}\xspace}
\newcommand{\bco}{\texttt{BCO}\xspace}
\newcommand{\oco}{\texttt{OCO}\xspace}
\newcommand{\hl}{\hat{\ell}\xspace}


\begin{document}
\maketitle
The goal of this document is design and analyze a simple algorithm for bandit convex optimization (BCO) in the continuous setting, which is based on the idea of optimistic estimators \citet{bubeck2016kernelbasedmethodsbanditconvex}.

\section{The complexity measure}
Instead of defining the whole bandit setting, regret, and the algorithm, we just focus on the complexity measure $\Lambda_q(E, p; f)$, which in turn results in a regret bound for the algorithm.

Given a convex body $\K$, a convex function $f:\K \to [0,1]$, two probability distributions $q$ and $p$ supported on $\K$, and an estimator $E:\K \times [0,1] \to \R^\cK$, we can define the complexity measure $\Lambda_q(E, p; f)$ as
\begin{align}
    \Lambda_q(E, p; f) := \int_{\K} p(t) f(t) dt - \int_{\K} q(t) \bar{f}(t) dt + \sup_{t\in \K} (\bar{f}(t) - f(t)) + \int_{\K} q(t) \int_{\K} p(x) E(x, f(x))_t^2 dx dt,
\end{align}
where $\bar{f}(t) := \int_{\K} p(x) E(x, f(x))_t dx$ is the estimator function associated with $E$ and $p$.
By using the notation $\langle f, p\rangle := \int_{\K} p(t) f(t) dt$, and letting $X\sim p$ and $T\sim q$ be random variables with distributions $p$ and $q$ respectively, we can rewrite $\Lambda_q(E, p; f)$ as
\begin{align}
    \Lambda_q(E, p; f) & = \langle f, p\rangle - \langle \bar{f}, q\rangle + \sup_{t\in \K} (\bar{f}(t) - f(t)) + \E_{T\sim q} \E_{X\sim p} [E(X, f(X))_T^2]                                                                                                                               \\
                       & = \underbrace{\langle f, p\rangle - \langle \bar{f}, q\rangle + \sup_{t\in \K} (\bar{f}(t) - f(t))}_{\triangleq \Lambda_{\mathrm{bias}, q}(E, p; f)} + \underbrace{\E_{T\sim q} \E_{X\sim p} [E(X, f(X))_T^2]}_{\triangleq \Lambda_{\mathrm{var}, q}(E, p; f)}\,.
\end{align}

\section{The kernel viewpoint}
To make the problem more tractable, we further restrict the search space of estimation functions.
Specifically, we only consider estimation functions of the form
\begin{align}
    E(x, y) = y \cdot s_x + b_x\,,
\end{align}
where $s_x, b_x \in \R^\K$ for $x \in \K$.
An affine estimator in this setting is a map
\begin{align}
    E: \cK \times[0,1]\to \R^\K,\qquad \text{which means} \qquad
    E(x,f(x))(\,\cdot\,)= f(x)\,s_x(\,\cdot\,) + b_x(\,\cdot\,),
\end{align}
and the average estimate is
\begin{equation}
    \bar f(t) := \int_\K p(x)\,E(x,f(x))(t)\,dx
    = \int_\K p(x)\,\bigl(f(x)s_x(t)+b_x(t)\bigr)\,dx.
    \label{eq:barf_def}
\end{equation}

\paragraph{The probabilistic view}
The overall goal is to design $S, B$ and $p$ jointly to minimize the worst-case value of $\Lambda_q$.
In this section we basically re-parametrize this joint design problem in a way that makes the structure of the solution more transparent.

Fix some $p$.
Define
\begin{align}
    Z(t) := \int_\K p(x) s_x(t) \,dx\,.
\end{align}
Suppose that $p(x)$ and $s_x(t)$ are such that $Z(t) \neq 0$ for all $t$ (which is the case for exponential kernels).
Then, we can define a new function $h(x|t)$ as
\begin{align}
    h(x|t) := \frac{p(x) s_x(t)}{Z(t)}\,, \qquad \text{so that} \qquad \int_\K h(x|t) \, dx = 1.
\end{align}

The idea is that \textbf{if we restrict to nonnegative kernels $s_x(t) \ge 0$},
then instead of designing $s_x(t)$ directly, we can equivalently design $h(\cdot|t)$ on $\K$ for each $t$,
and the scaling factors $Z_t$ for each $t$.

Starting from the continuous form $\bar{f}$ , we have
\begin{align}
    \bar{f}(t) & = \int_{\K} p(x) s_x(t) f(x) \,dx + \underbrace{\int_{\K} p(x)\,b_x(t)}_{:=C(t)}\,dx \\
               & = Z(t) \cdot \int_{\K} h(x|t) f(x) \, dx + C(t)                                      \\
               & = Z(t) \cdot \E_{x\sim h(\cdot|t)}[f(x)] + C(t)\,.
\end{align}
The take away message is that, when we restrict to nonnegative kernels, the design of the estimator can be viewed as the design
of a family of (kernel) distributions $h(\cdot|t)$ on $\K$ for each $t$, and the scaling and bias factors $Z(t)$ and $C(t)$ for each $t$.

\section{The optimistic estimator}
We switch gear and consider a new idea for designing estimators based on \cite{bubeck2016kernelbasedmethodsbanditconvex}, this time in the continuous setting
in $\R^d$.
Let $\K \subset \Rd$ be a convex body, and $f: \K \to [0,1]$ be a convex function supported on $K$.
We define
\begin{align}
    \tilde{f}(t) = \E\left[
        \left(1 - \frac{1}{\lambda}\right) f(X)
        + \frac{1}{\lambda} f\left(
        (1-\lambda) X + \lambda t
        \right)
        \right]\,,
\end{align}
where $X \sim p$ is a random variable supported on $\K$ and $\lambda\in(0,1)$ is a parameter.

First, it is not clear if $\tilde{f}(t)$ is a good estimate of $f(t)$;
however, even if it is, \textit{it is not a valid $\bar{f}$},
since we do not yet know if there exists an estimator function $E$ such that
\begin{align}
    \tilde{f}(t) \stackrel{?}{=} \int_{\Rd} p(x) E(x,f(x))(t)\,dx.
\end{align}
But before we worry about the existence of an estimator, let's first check if $\tilde{f}(t)$ is a good estimate of $f(t)$.
The key observation is that \textbf{every term in the expectation defining $\tilde{f}(t)$ is an underestimate of $f(t)$, hence $\tilde{f}(t)\le f(t)$ for all $t$}.
To see that, note that by convexity we have
\begin{align}
    f\left((1-\lambda)X+\lambda t\right)
    \le (1-\lambda)f(X) + \lambda f(t),
\end{align}so
\begin{align}
    \left(1 - \frac{1}{\lambda}\right) f(X)
    + \frac{1}{\lambda} f\left((1-\lambda) X + \lambda t\right)
    \le f(t).
\end{align}
Therefore, $\tilde{f}(t)$ is a pointwise underestimate of $f(t)$, i.e
\begin{align}
    \label{eq:tilde_f_underestimates_f}
    \tilde{f}(t) \le f(t),\qquad \forall t\in \Rd.
\end{align}
\paragraph{The valid estimator $E$ for $\tilde{f}$.}
With a change of measure argument, we can show that there exists an estimator $E$ for $\tilde{f}$.
Precisely, we can write
\begin{align}
    \tilde{f}(t)
     & = \int_{\Rd}
    p(x) \left[
        \left(1 - \frac{1}{\lambda}\right) f(x)
        + \frac{1}{\lambda} f\left(
        (1-\lambda) x + \lambda t
        \right)
    \right]\,dx     \\
     & =
    \left(1 - \frac{1}{\lambda}\right)
    \int_{\Rd}
    p(x)
    f(x)
    dx
    +
    \frac{1}{\lambda}
    \int_{\Rd}
    p(x)
    f\left(
    (1-\lambda) x + \lambda t
    \right) \,dx    \\
     & =
    \left(1 - \frac{1}{\lambda}\right)
    \int_{\Rd}
    p(x)
    f(x)
    dx
    +
    \frac{1}{\lambda (1-\lambda)^d}
    \int_{\Rd}
    p\left(
    \frac{y - \lambda t}{1 - \lambda}
    \right)
    f(y) \,dy       \\
     & =
    \int_{\Rd}
    p(x)
    f(x)
    \left(
    1 - \frac{1}{\lambda}
    +
    \frac{1}{\lambda (1-\lambda)^d}
    \cdot
    p\left(
    \frac{x - \lambda t}{1 - \lambda}
    \right)
    p(x)^{-1}
    \right) \,dx\,,
\end{align}
where we have used the change of variable $y = (1-\lambda)x + \lambda t$ in the second term.
Therefore, we can define the estimator $E$ as
\begin{align}
    E_\lambda(x, f(x))_t := f(x) \cdot \left(1 - \frac{1}{\lambda} + \frac{1}{\lambda (1-\lambda)^d} \cdot p\left(\frac{x - \lambda t}{1 - \lambda}\right) p(x)^{-1}\right)\,.
\end{align}
An interesting observation is that $E_\lambda(x, f(x))_t$ is a linear (more than affine) function of $f(x)$,
so using our notation for affine estimators, we can write $E_\lambda(x, f(x))_t = f(x) \cdot s^\lambda_x(t)$, where
\begin{align}
    s^\lambda_x(t) := 1 - \frac{1}{\lambda} + \frac{1}{\lambda (1-\lambda)^d} \cdot p\left(\frac{x - \lambda t}{1 - \lambda}\right) p(x)^{-1}\,.
\end{align}
Further, notice that the density of the random variable $Y = (1-\lambda)X + \lambda t$ is given by
\begin{align}
    p^\lambda_t(y) = \frac{1}{(1-\lambda)^d} p\left(\frac{y - \lambda t}{1 - \lambda}\right)\,,
\end{align}
and then we have
\begin{align}
    p(x) s^\lambda_x(t)
     & = p(x) \left(1 - \frac{1}{\lambda}\right) + \frac{1}{\lambda} p^\lambda_t(x) \,,
\end{align}
which means that $p(x) s^\lambda_x(t)$ is an affine combination of two probability densities $p$ and $p^\lambda_t$,
and therefore the following normalization constant is exactly 1:
\begin{align}
    Z(t) & = \int_{\Rd} p(x) s^\lambda_x(t) \,dx                                                                                                                       \\
         & = \int_{\Rd} p(x) \left(1 - \frac{1}{\lambda} + \frac{1}{\lambda (1-\lambda)^d} \cdot p\left(\frac{x - \lambda t}{1 - \lambda}\right) p(x)^{-1}\right) \,dx \\
         & =
    1 - \frac{1}{\lambda} + \frac{1}{\lambda (1-\lambda)^d} \int_{\Rd} p\left(\frac{x - \lambda t}{1 - \lambda}\right) \,dx                                            \\
         & = 1.
\end{align}
So, now we can see what $h(\cdot|t)$ is in this case:
\begin{align}
    h_\lambda(x|t) := \frac{p(x) s^\lambda_x(t)}{Z(t)}
     & = p(x) \left(1 - \frac{1}{\lambda}\right) + \frac{1}{\lambda} p^\lambda_t(x) \,.
\end{align}
Let's see what is $\E_{X\sim h_\lambda(\cdot|t)}[X]$ in this case:
\begin{align}
    \E_{X\sim h_\lambda(\cdot|t)}[X]
     & = \int_{\Rd} h_\lambda(x|t) x \,dx
    \\
     & = \left(1 - \frac{1}{\lambda}\right) \int_{\Rd} p(x) x \,dx + \frac{1}{\lambda} \int_{\Rd} p_t^\lambda\left(x\right)  x \,dx
    \\
     & =
    \left(1 - \frac{1}{\lambda}\right)
    \int_{\Rd} p(x)  x \,dx
    + \frac{1}{\lambda}
    \int_{\Rd} p\left(y\right) \left(
    (1-\lambda) y + \lambda t
    \right) \,dy
    \\
     & = \left(
    1 - \frac{1}{\lambda}
    \right) \E[X] + \frac{1}{\lambda} \left(
    (1-\lambda) \E[X] + \lambda t
    \right)
    = t,
\end{align}
therefore, we have the barycenter condition \eqref{eq:barycenter_condition} satisfied in this case.
\begin{remark}
    It is tempting to apply Jensen's inequality to $f$ and get
    \begin{align}
        \bar{f}(t)
        = \E_{X\sim h_\lambda(\cdot|t)}[f(X)]
        \stackrel{(\text{Jensen})}{\geq}
        f\left(\E_{X\sim h_\lambda(\cdot|t)}[X]\right)
        = f(t),
    \end{align}
    so by combining this with \eqref{eq:tilde_f_underestimates_f}, do we (almost) have a contradiction?
    No.
    The issue is that $s^\lambda_x(t)$ can be negative for some $x$ and therefore,
    Jensen's inequality does not apply to $s(t) = \int_{\Rd} p(x) s^\lambda_x(t) f(x) dx$.
    In fact, $h_\lambda(\cdot|t)$ is not a probability distribution,
    rather it is a signed measure, which integrates to 1 but can have negative mass on some sets.
\end{remark}

\paragraph{Taking a look at $\Lambda_{\mathrm{bias}}$}
First, we need to actually define the $\Lambda_{\mathrm{bias}}$ term for the continuous setting, which is
\begin{align}
    \Lambda_{\mathrm{bias}, q}(E_\lambda, p; f)
     & =
    \int_{\Rd} p(t) f(t) dt - \int_{\Rd} q(t) \bar{f}(t) dt + \sup_{t\in \K} (\bar{f}(t) - f(t))\,.
\end{align}
We need to somehow \textit{bind} $p$ to $q$.
This time, since $\bar{f}$ is a pointwise underestimate of $f$, we would be in good shape
if we can do our kernel trick to make $\int_{\Rd} q(t) \bar{f}(t) dt = \int_{\Rd} p(t) f(t) dt$.
So to this end, we can define $p$ as
\begin{align}
    p(x) := \int_{\Rd} q(t) h_\lambda(x|t) dt\,.
\end{align}
Our concern should be whether $p$ is a valid probability distribution.
One can check that $\int_{\Rd} p(x) dx = 1$, so we only need to check if $p(x)\ge 0$ for all $x$.
To this end, we have
\begin{align}
    p(x) & = \int_{\Rd} q(t) h_\lambda(x|t) dt                                                                          \\
         & = \int_{\Rd} q(t) \left(p(x) \left(1 - \frac{1}{\lambda}\right) + \frac{1}{\lambda}p_t^\lambda(x) \right) dt \\
         & = \left(1 - \frac{1}{\lambda}\right) p(x)
    +
    \frac{1}{\lambda} \int_{\Rd} q(t) p_t^\lambda\left(x\right) dt\,,
\end{align}
so technically for $p$ to be even well-defined, we need
\begin{align}
    p(x) = (1 - \lambda)^{-d} \int_{\Rd} q(t) p\left(\frac{x - \lambda t}{1 - \lambda}\right) dt\,,
\end{align}
which is a convolution-type equation for $p$.
To analyze this, define the operator $\cT$ on the space of probability densities on $\Rd$ as
\begin{align}
    (\cT p)(x) := (1 - \lambda)^{-d} \int_{\Rd} q(t) p\left(\frac{x - \lambda t}{1 - \lambda}\right) dt\,,
\end{align}
which is exactly the density of the random variable
\begin{align}
    X' = (1-\lambda) X + \lambda T\,,
\end{align}
where $X\sim p$ and $T\sim q$ are independent.
Therefore, the condition for $p$ to be well-defined is exactly that $p$ is a fixed point of $\cT$, i.e.
\begin{equation}
    p = \cT p\,.
    \label{eq:fixed_point_condition}
\end{equation}
which means that $p$ is a stationary distribution for the affine recursion (or Markov chain) of
\begin{align}
    X_{n+1} = (1-\lambda) X_n + \lambda T_{n+1}, \qquad
    T_n \stackrel{iid}{\sim} q\,.
\end{align}
Define
\begin{align}
    X := \lambda \sum_{k=0}^\infty (1-\lambda)^k T_{k},
\end{align}
which converges almost surely whenever $q$ has a finite first moment.
Then, the law of $X$ is a fixed point of $\cT$, since
\begin{align}
    (1-\lambda) X + \lambda T_0
     & = (1-\lambda) \lambda \sum_{k=0}^\infty (1-\lambda)^k T_{k} + \lambda T_0
    = \lambda \sum_{k=0}^\infty (1-\lambda)^k T_{k} = X.
\end{align}
Therefore, the law of $X$ is a valid choice for $p$,
and it is the unique stationary solution.
Also, one can show that $p$ is supported on $\cK$ if $q$ is supported on $\cK$.
This is a very big win for the continuous setting, since we now have
\begin{align}
    \Lambda_{\mathrm{bias}, q}(E_\lambda, p; f)
     & =
    \int_{\Rd} p(t) f(t) dt - \int_{\Rd} q(t) \bar{f}(t) dt + \sup_{t\in \K} (\bar{f}(t) - f(t))
    \\
     & =
    \sup_{t\in \K} (\bar{f}(t) - f(t))
    \le 0\,,
\end{align}
which is surprisingly non-positive.
\begin{remark}
    Some immediate moment relations between $X \sim p$ and $T\sim q$:
    \begin{itemize}
        \item We have
              \begin{align}
                  \E[X] & = \lambda \sum_{k=0}^\infty (1-\lambda)^k \E[T_k] = \E[T]\,.
              \end{align}
        \item We have
              \begin{align}
                  \mathrm{Cov}[X]
                   & = \lambda^2 \sum_{k=0}^\infty (1-\lambda)^{2k} \mathrm{Cov}[T_k]
                  = \frac{\lambda}{2-\lambda} \mathrm{Cov}[T].
              \end{align}
    \end{itemize}
\end{remark}

\begin{example}
    If $q$ is $\cN(\mu, \Sigma)$, then $p$ is $\cN(\mu, \frac{\lambda}{2-\lambda}\Sigma)$.
\end{example}

\paragraph{The variance term.}
The variance term  in the continuous setting is
\begin{align}
    \Lambda_{\text{var}, q}(E_\lambda, p; f)
     & =
    \int_{\Rd} q(t) \int_{\Rd} p(x) E_\lambda(x, f(x))_t^2 dx dt \\
     & = \int_{\Rd} p(x) q(t) f(x)^2 s^\lambda_x(t)^2 dt dx\,.
\end{align}

Let's once more upper bound $f(x)$ by 1 (not too naively) to get
\begin{align}
    \Lambda_{\text{var}, q}(E_\lambda, p; f)
    \leq
    \Lambda_{\text{var}, q}(E_\lambda, p; \mathbf{1})
     & \le \int_{\Rd} p(x) q(t) s^\lambda_x(t)^2 dt dx\,,
\end{align}
which might not be as bad as it looks, since this is actually the worst case for the variance
term, which is realized by the constant function $f\equiv 1$.
Recall that
\begin{align}
    p_t^\lambda(x) := \frac{1}{(1-\lambda)^{d}}p\left(
    \frac{
        x - \lambda t
    }{
        1 - \lambda
    }
    \right)\,,
\end{align}
which is the density of the random variable $Y\triangleq(1-\lambda)X + \lambda t$ when $X\sim p$.
Let,
\begin{align}
    \omega_t(x) := \frac{p_t^\lambda(x)}{p(x)}\,,
\end{align}
then for every $t$ we have
\begin{align}
    \int_{\Rd} p(x) \omega_t(x) dx
     & = \int_{\Rd} p_t^\lambda(x) dx
    = 1\,,
\end{align}
so $\omega_t$ is a density ratio function with respect to $p$.
Now, observe that $\int (\omega_t(x) - 1) p(x) dx = 0$, so $\omega_t(x)-1$ is a zero-mean function with respect to $p$,
and recalling that $p(x)s^\lambda_x(t) = (1 - \frac{1}{\lambda})p(x) + \frac{1}{\lambda} \cdot p_t^\lambda\left(x\right)$,
we can write
\begin{align}
    \int p(x) s_x^\lambda(t)^2 dx
     & = \int p(x) \left(1 + \frac{1}{\lambda}(\omega_t(x) - 1)\right)^2 dx   \\
     & = \int p(x) \left(1 + \frac{1}{\lambda^2}(\omega_t(x) - 1)^2\right) dx \\
     & = 1 + \frac{1}{\lambda^2}\int p(x) (\omega_t(x) - 1)^2 dx\,,
\end{align}
but using the definition of $\omega_t$ we can also write
\begin{align}
    1 + \frac{1}{\lambda^2}\int p(x) (\omega_t(x) - 1)^2 dx\,,
     & = 1 + \frac{1}{\lambda^2}\int p(x) (\frac{p_t^\lambda(x)}{p(x)} - 1)^2 dx \\
     & = 1 + \frac{1}{\lambda^2} \chi^2(p_t^\lambda \| p)\,,
\end{align}
where $\chi^2(p_t^\lambda \| p)$ is the chi-squared divergence between $p_t^\lambda$ and $p$.
Therefore, in one line we have
\begin{align}
    \int p(x) s_x^\lambda(t)^2 dx
     & = 1 + \frac{1}{\lambda^2} \chi^2(p_t^\lambda \| p)\,,
\end{align}
and consequently, we have
\begin{align}
    \sup_{f \in [0,1], \textrm{convex}}\Lambda_{\text{var}, q}(E_\lambda, p; f)
     & \stackrel{(\mathrm{"="for} f\equiv 1)}{\le} \int q(t) \left(1 + \frac{1}{\lambda^2} \chi^2(p_t^\lambda \| p)\right) dt \\
     & = 1 + \frac{1}{\lambda^2} \E_{T\sim q}[\chi^2(p_T^\lambda \| p)]\,.
\end{align}
Recall the setup
\begin{align}
    T \sim q,\,  \qquad \qquad X \sim p, \qquad \qquad
    Y                = (1-\lambda) X + \lambda T\, ,
\end{align}
and $p_t^\lambda$ is the density of $Y|T=t$.
Now by expanding the chi-squared divergence, we can write
\begin{align}
    \E_{T\sim q}[\chi^2(p_T^\lambda \| p)]
     & =
    \int \int
    \frac{
        (p_t^\lambda(y) - p(y))^2
    }{p(y)} q(t) \, dy dt \\
     & =
    \int \int
    \frac{
        \left(q(t)p_t^\lambda(y) - q(t)p(y)\right)^2
    }{q(t) p(y)} \, dy dt\,,
\end{align}
and note that $Y \stackrel{D.}{=} X$, so we can write
\begin{align}
    \int \int
    \frac{
        \left(q(t)p_t^\lambda(y) - q(t)p(y)\right)^2
    }{q(t) p(y)} \, dy dt
     & =
    \int \int
    \frac{
        \left(P_{Y,T}(y,t) - P_T(t)P_X(y)\right)^2
    }{P_X(y) P_T(t)} \, dy dt                            \\
     & =
    \int \int
    \frac{
        \left(P_{Y,T}(y,t) - P_T(t)P_Y(y)\right)^2
    }{P_Y(y) P_T(t)} \, dy dt                            \\
     & = \chi^2(P_{Y,T} \| P_Y P_T) = I_{\chi^2}(T;Y)\,,
\end{align}
where $I_{\chi^2}(T;Y)$ is the chi-squared mutual information between $T$ and $Y$.
Note that generally speaking, by decreasing $\lambda$ the mutual information $I_{\chi^2}(T;Y)$ decreases, but the factor $\frac{1}{\lambda^2}$ increases, so there is a tradeoff here.

\begin{example}[The one dimensional Gaussian case]
    Let $q$ be $\cN(\mu, \sigma_q^2)$, then $p$ is $\cN(\mu, \sigma_p^2)$,
    where $\sigma_p^2 = \frac{\lambda}{2-\lambda} \sigma_q^2$.
    Further, $\mu_t$ is the density of $Y = (1 - \lambda) X + \lambda T$,
    where $X \sim p$ and $T \sim q$ are independent.
    The pair $(T,Y)$ is jointly Gaussian with corelation
    \begin{align}
        \rho^2
         & = \frac{\mathrm{Cov}[T,Y]^2}{\mathrm{Var}[T]\mathrm{Var}[Y]}
        = \frac{\mathrm{Cov}[T, (1-\lambda) X + \lambda T]^2}{\mathrm{Var}[T]\mathrm{Var}[(1-\lambda)X + \lambda T]} \\
         & = \frac{\left((1-\lambda) \mathrm{Cov}[T, X] + \lambda \mathrm{Cov}[T, T] \right)^2}{
            \mathrm{Var}[T]
            \left(
            (1-\lambda)^2 \mathrm{Var}[X] + \lambda^2 \mathrm{Var}[T]
            \right)
        }
        = \frac{\lambda^2 \mathrm{Var}[T]^2}{
            \mathrm{Var}[T]
            \left(
            (1-\lambda)^2 \mathrm{Var}[X] + \lambda^2 \mathrm{Var}[T]
            \right)
        }                                                                                                            \\
         & = \frac{\lambda^2\sigma_q^2}{
            (1-\lambda)^2 \sigma_p^2 + \lambda^2 \sigma_q^2
        }
        = \frac{\lambda^2}{
            \frac{\lambda(1-\lambda)^2}{2-\lambda} + \lambda^2
        } = \lambda(2 - \lambda)\,,
    \end{align}
    and therefore $(1 - \rho^2) = (1-\lambda)^2$.
    Now, using the formula for the chi-squared mutual information between jointly Gaussian variables, we have
    \begin{align}
        \chi^2(P_{X,T} \| P_X P_T) = I_{\chi^2}(T;X) = \frac{\rho^2}{1-\rho^2} = \frac{\lambda(2-\lambda)}{(1-\lambda)^2}
        =(1 - \lambda)^{-2} - 1\,.
    \end{align}
    Hence, we have
    \begin{align}
        \sup_{f \in [0,1], \textrm{convex}}\Lambda_{\text{var}, q}(E_\lambda, p; f)
         & = 1 + \frac{1}{\lambda^2} \left(
        (1 - \lambda)^2 - 1
        \right)
        = 1 + \frac{2 - \lambda}{\lambda (1-\lambda)^2}\,.
    \end{align}
\end{example}

\begin{example}[The $d$-dimensional Gaussian case]
    If $q$ is $\cN(\mu, \Sigma_q)$, then $p$ is $\cN(\mu, \Sigma_p)$,
    where $\Sigma_p = \frac{\lambda}{2-\lambda} \Sigma_q$.
    In whitened coordinates, the $(X, T)$ joint splits into $d$ independent
    correlated Gaussian pairs with the same $\rho^2 = \lambda(2-\lambda)$.
    Therefore, we have
    \begin{align}
        I_{\chi^2}(T;X) = \prod_{i=1}^d \frac{1}{1-\rho^2} - 1 = (1-\lambda)^{-2d} - 1\,,
    \end{align}
    and
    \begin{align}
        \sup_{f \in [0,1], \textrm{convex}}\Lambda_{\text{var}, q}(E_\lambda, p; f)
         & = 1 + \frac{1}{\lambda^2} \left(
        (1 - \lambda)^{-2d} - 1
        \right)\,.
    \end{align}
    Now by setting $\lambda = \frac{c}{d}$ for some constant $c>0$, we get
    \begin{align}
        (1 - \lambda)^{-2d} - 1
         & = \left(1 - \frac{c}{d}\right)^{-2d} - 1
        \approx e^{2c} - 1\,,
    \end{align}
    so we have
    \begin{align}
        \sup_{f \in [0,1], \textrm{convex}}\Lambda_{\text{var}, q}(E_\lambda, p; f)
         & \approx 1 + \frac{e^{2c} - 1}{c^2} d^2\,,
    \end{align}
    which is a quadratic dependence on $d$.
\end{example}

So the bottom line is
\begin{align}
    \sup_{f \in [0,1], \textrm{convex}}\Lambda_{\text{var}, q}(E_\lambda, p; f)
     & = 1 + \frac{1}{\lambda^2} I_{\chi^2}(T;Y)\,,
\end{align}
where $T\sim q$, $X\sim p$ and $Y = (1-\lambda) X + \lambda T$.
However, recall that
\begin{align}
    X = \lambda \sum_{k=1}^\infty (1-\lambda)^{k-1} T_k\,,
\end{align}
where $T_k$ are i.i.d. copies of $T$, so $Y$ can be written as
\begin{align}
    Y = (1-\lambda) X + \lambda T = \lambda \sum_{k=1}^\infty (1-\lambda)^{k} T_k + \lambda T_0
    = \lambda \sum_{k=0}^\infty (1-\lambda)^k T_k\,,
\end{align}
where we set $T_0 \triangleq T$ in the second equality.
Therefore, we can equally write
\begin{equation}
    I_{\chi^2}(T;Y) = I_{\chi^2}\left(T_0; \lambda \sum_{k=0}^\infty (1-\lambda)^k T_k\right)\,.
    \label{eq:chi_squared_mutual_info_infinite_sum}
\end{equation}
To step a bit further from Gaussians, we can ask what would happen if $q$ is a log-concave distribution
on a convex body $\cK$; then what would be the best bound one can achieve on \eqref{eq:chi_squared_mutual_info_infinite_sum}?

\clearpage
\subsection{Bounding the chi-squared mutual information}
We are looking for a set of regularity conditions on $q$ that would give us a good upper bound on the chi-squared mutual information quantity in the following:
\begin{conjecture}\label{conj:logconcave_chi2_poly_d}
    Let $q$ be a log-concave distribution on a convex body $\cK\subseteq \Rd$, and let $T_k$ be i.i.d. samples from $q$.
    Does there exist a $\poly(d)$ upper bound on the following chi-squared mutual information quantity?
    \begin{align}
        \sup_{q} \inf_{\lambda\in(0,1)} \frac{1}{\lambda^2} I_{\chi^2}\left(T_0; \lambda \sum_{k=0}^\infty (1-\lambda)^k T_k\right)\,.
    \end{align}
\end{conjecture}

\subsubsection{The mutual information is infinite for Laplace}
Suppose that $T \sim \mathrm{Laplace}(0, 1)$, then
\begin{align}
    q(t) = \frac{1}{2} e^{-|t|}\,,
\end{align}
and $\lambda T$ has density $q_\lambda(t) \triangleq \frac{1}{2\lambda} e^{-\frac{|t|}{\lambda}}$.
Recall that $Y = (1-\lambda) X + \lambda T$, where $X$ is an independent copy of $Y$, so $p = q_\lambda * p$, and
for every $y$ we have
\begin{align}
    p(y) = \int q_\lambda(y - z) r(z) dz = \E_{Z \sim r} \left[q_\lambda(y - Z)\right]
    = \E_{Z \sim r} \left[\frac{1}{2\lambda} e^{-\frac{|y-Z|}{\lambda}}\right]\,,
\end{align}
For $y > 0$, we have the triangle inequality $|y-Z| \ge y - |Z|$, so we can write
\begin{align}
    p(y) \le \E_{Z \sim r} \left[\frac{1}{2\lambda} e^{\frac{-y + |Z|}{\lambda}}\right] = \frac{1}{2\lambda} e^{-\frac{y}{\lambda}} \E_{Z \sim r} \left[e^{\frac{|Z|}{\lambda}}\right]\,,
\end{align}
which by setting $C_\lambda \triangleq \E_{Z \sim r} \left[e^{\frac{|Z|}{\lambda}}\right]$ gives
\begin{align}
    p(y) \le \frac{C_\lambda}{2\lambda} e^{-\frac{y}{\lambda}}\,.
\end{align}
However, we need to show that $C_\lambda$ is finite. To this end, write
\begin{align}
    Z = \lambda \sum_{k=1}^\infty (1-\lambda)^k T_k \quad \Rightarrow \quad \frac{|Z|}{\lambda} \leq \sum_{k=1}^\infty (1-\lambda)^k |T_k|
    \quad \Rightarrow \quad \exp(\frac{|Z|}{\lambda}) \leq \prod_{k=1}^\infty \exp((1-\lambda)^k |T_k|)\,,
\end{align}
then taking expectation and using the independence of $T_k$ gives
\begin{align}
    \E_{Z \sim r} \left[e^{\frac{|Z|}{\lambda}}\right] \leq \prod_{k=1}^\infty \E_{T_k \sim q} \left[e^{(1-\lambda)^k |T_k|}\right]\,,
\end{align}
but for a laplace variable $L$ and $\alpha \in (0,1)$ we have
\begin{align}
    \E[e^{\alpha |L|}] = \int_0^\infty e^{-x} e^{\alpha x} dx = \int_0^\infty e^{x(\alpha - 1)}dx = \frac{1}{1 - \alpha}\,.
\end{align}
Therefore, we have
\begin{align}
    \E_{Z \sim r} \left[e^{\frac{|Z|}{\lambda}}\right] \leq \prod_{k=1}^\infty \frac{1}{1 - (1-\lambda)^k}\,,
\end{align}
To show that this is finite, it is enough to show that
\begin{align}
    \sum_{k=1}^\infty -\log\bigl(1-(1-\lambda)^k\bigr) < \infty\,.
\end{align}
Since $(1-\lambda)^k \to 0$, there exists $K$ such that for all $k \ge K$,
\begin{align}
    (1-\lambda)^k \le \frac12\,.
\end{align}
Now for every $x \in [0,1/2]$ we have
\begin{align}
    -\log(1-x) \le 2x\,.
\end{align}
Therefore, for all $k \ge K$,
\begin{align}
    -\log\bigl(1-(1-\lambda)^k\bigr)
    \le 2(1-\lambda)^k\,.
\end{align}
Summing over $k \ge K$ gives
\begin{align}
    \sum_{k=K}^\infty -\log\bigl(1-(1-\lambda)^k\bigr)
    \le 2\sum_{k=K}^\infty (1-\lambda)^k < \infty\,,
\end{align}
since the right-hand side is a convergent geometric series. The finitely many terms with $k < K$ are also finite, so altogether
\begin{align}
    \sum_{k=1}^\infty -\log\bigl(1-(1-\lambda)^k\bigr) < \infty\,.
\end{align}
Hence
\begin{align}
    \prod_{k=1}^\infty \frac{1}{1-(1-\lambda)^k} < \infty\,,
\end{align}
and therefore $C_\lambda < \infty$. We conclude that for all $y>0$,
\begin{align}
    p(y) \le \frac{C_\lambda}{2\lambda} e^{-y/\lambda}\,.
\end{align}

Next, note that
\begin{align}
    Z = \lambda(1-\lambda)T_1 + \lambda\sum_{k=2}^\infty (1-\lambda)^k T_k\,.
\end{align}
Thus the law $r$ of $Z$ is the convolution of a Laplace density with a probability measure, and in particular $r$ is continuous and strictly positive everywhere. Hence there exist $\epsilon>0$ and $m_\lambda>0$ such that
\begin{align}
    r(u) \ge m_\lambda
    \qquad\text{for all } |u|\le \epsilon\,.
\end{align}

Now recall that the conditional law of $Y$ given $T=t$ is exactly the translate of $r$:
\begin{align}
    p_t^\lambda(y)=r(y-\lambda t)\,.
\end{align}
Therefore, by the integral formula for $\chi^2$-mutual information,
\begin{align}
    I_{\chi^2}(T;Y)+1
     & = \int_{\mathbb R}\int_{\mathbb R}
    q(t)\,\frac{p_t^\lambda(y)^2}{p(y)}\,dy\,dt \\
     & = \int_{\mathbb R}\int_{\mathbb R}
    q(t)\,\frac{r(y-\lambda t)^2}{p(y)}\,dy\,dt\,.
\end{align}
Restrict the integral to the strip
\begin{align}
    \mathcal S \triangleq \{(t,y): t\ge 2\epsilon/\lambda,\ |y-\lambda t|\le \epsilon\}\,.
\end{align}
On this strip we have $y\ge \epsilon>0$, so the tail bound on $p$ applies, and also
\begin{align}
    r(y-\lambda t)\ge m_\lambda\,.
\end{align}
Hence for $(t,y)\in\mathcal S$,
\begin{align}
    q(t)\frac{r(y-\lambda t)^2}{p(y)}
     & \ge \frac12 e^{-t}\,
    \frac{m_\lambda^2}{(C_\lambda/(2\lambda))e^{-y/\lambda}} \\
     & = \frac{\lambda m_\lambda^2}{C_\lambda}
    \exp\!\left(-t+\frac{y}{\lambda}\right).
\end{align}
But on $\mathcal S$ we have $y=\lambda t + u$ with $|u|\le \epsilon$, so
\begin{align}
    -t+\frac{y}{\lambda}
    = -t + \frac{\lambda t + u}{\lambda}
    = \frac{u}{\lambda}
    \ge -\frac{\epsilon}{\lambda}.
\end{align}
Therefore, throughout $\mathcal S$,
\begin{align}
    q(t)\frac{r(y-\lambda t)^2}{p(y)}
    \ge \frac{\lambda m_\lambda^2}{C_\lambda}e^{-\epsilon/\lambda}
    =: c_\lambda > 0\,.
\end{align}
It follows that
\begin{align}
    I_{\chi^2}(T;Y)+1
    \ge \iint_{\mathcal S} c_\lambda\,dy\,dt.
\end{align}
Since $\mathcal S$ has infinite Lebesgue measure in the $t$-direction, the right-hand side is infinite. Hence
\begin{align}
    I_{\chi^2}(T;Y)=\infty.
\end{align}
Since this holds for every $\lambda\in(0,1)$, we conclude that
\begin{align}
    \inf_{\lambda\in(0,1)} \frac{1}{\lambda^2} I_{\chi^2}(T;Y)=\infty.
\end{align}

\clearpage
\subsubsection{The KL mutual information case.}
It is instructive to compare Conjecture \ref{conj:logconcave_chi2_poly_d} with the usual (KL-based) mutual information.
Assume $q$ has a density on a full-dimensional convex body $\cK$ (which holds for log-concave measures on $\cK$),
so that $Y$ also has a density and finite entropy.
Recall the decomposition
\begin{align}
    Y
    = \lambda T_0 + (1-\lambda)Y',
    \qquad
    Y' \triangleq \lambda \sum_{k=0}^\infty (1-\lambda)^k T_{k+1},
\end{align}
where $Y'$ is independent of $T_0$ and $Y' \stackrel{d}{=} Y$.
Let $h(\cdot)$ denote differential entropy.
Then the KL mutual information admits an exact, distribution-free evaluation:
\begin{align}
    I(T_0;Y)
     & = h(Y) - h(Y \mid T_0)\label{eq:kl_mi_entropy_def}                                                                     \\
     & = h(Y) - \E_{T_0}\!\left[ h\bigl( (1-\lambda)Y' + \lambda T_0 \,\big|\, T_0 \bigr) \right]\label{eq:kl_mi_expand_cond} \\
     & = h(Y) - h\bigl((1-\lambda)Y'\bigr)\label{eq:kl_mi_translation_invariance}                                             \\
     & = h(Y) - \bigl(h(Y') + d\log(1-\lambda)\bigr)\label{eq:kl_mi_scaling}                                                  \\
     & = d\log\!\left(\frac{1}{1 - \lambda}\right)\label{eq:kl_mi_final},
\end{align}
where \eqref{eq:kl_mi_translation_invariance} uses translation invariance of differential entropy
(the shift by $\lambda T_0$ does not affect $h(\cdot)$), \eqref{eq:kl_mi_scaling} uses the scaling rule
$h(aX)=h(X)+d\log|a|$, and \eqref{eq:kl_mi_final} uses $Y'\stackrel{d}{=}Y$.

In particular, unlike the $\chi^2$-mutual information, the quantity $I(T_0;Y)$ depends only on $\lambda$ and $d$,
and is completely independent of the choice of $q$:
\begin{align}
    I(T_0;Y) = d\log\!\left(\frac{1}{1-\lambda}\right).
\end{align}
Consequently, $\inf_{\lambda\in(0,1)} I(T_0;Y)=0$ (by taking $\lambda\downarrow 0$), while for the scaling
$\lambda = c/d$ one obtains the constant-order behavior $I(T_0;Y)=c+O(1/d)$.

\clearpage

\subsubsection{An upper bound based on $\|q\|_{\infty}$}
Let $Y := \lambda \sum_{k=0}^\infty (1-\lambda)^k T_k$ be the infinite autoregressive sum,
then by the definition of the chi-squared divergence, we can write
\begin{align}
    I_{\chi^2}(T_0; Y)
     & = \chi^2(P_{T_0, Y} \| P_{T_0} P_{Y})                                                       \\
     & = \int_{\cK} \int_{\cK} \frac{P_{T_0, Y}(t, y)^2}{P_{T_0}(t) P_{Y}(y)} dy dt - 1            \\
     & = \int_{\cK} \int_{\cK} \frac{P_{T_0}(t)^2 P_{Y|T_0=t}(y)^2}{P_{T_0}(t) P_{Y}(y)} dy dt - 1 \\
     & = \int_{\cK} \int_{\cK} \frac{P_{T_0}(t) P_{Y|T_0=t}(y)^2}{P_{Y}(y)} dy dt - 1              \\
     & = \int_{\cK} P_{T_0}(t) \left(\int_{\cK} \frac{P_{Y|T_0=t}(y)^2}{P_{Y}(y)} dy - 1\right) dt \\
     & = \E_{T_0} \left[
        \chi^2(P_{Y|T_0} \| P_{Y})
        \right]\,.
\end{align}
Note that
\begin{align}
    Y = \lambda T_0 + \lambda \sum_{k=1}^\infty (1-\lambda)^k T_k
    = \lambda T_0 + (1-\lambda) \lambda \sum_{k=0}^\infty (1-\lambda)^k T_{k+1}
    = \lambda T_0 + \underbrace{(1-\lambda) Y^\prime}_{\triangleq Z}\,,
\end{align}
where $Y^\prime$ is an independent copy of $Y$.
Define the densities of the relevant distributions as
\begin{align}
    q(t) \triangleq P_{T_0}(t), \quad p(y) \triangleq P_{Y}(y), \quad \text{and} \quad p_t^\lambda(y) \triangleq P_{Y|T_0=t}(y)\,, \quad \text{and} \quad r(z) \triangleq P_Z(z)\,.
\end{align}
Although the exact form of $p$ in terms of $q$ is not very explicit,
the other two densities are just scaled shifts of $p$.
First, note that $Z = (1-\lambda) Y^\prime$, so we have
\begin{align}
    r(z) = \frac{1}{(1-\lambda)^d} p\left(\frac{z}{1-\lambda}\right)\,,
\end{align}
and secondly, given $T_0 = t$, we have $Y = \lambda t + Z$, so the density of $Y$ given $T_0 = t$ is just a shifted version of the density of $Z$:
\begin{align}
    p_t^\lambda(y) = r(y - \lambda t) = \frac{1}{(1-\lambda)^d} p\left(\frac{y - \lambda t}{1-\lambda}\right)\,.
\end{align}
Now, by putting these back into the chi-squared divergence we have
\begin{align}
    \chi^2(P_{Y|T_0=t} \| P_{Y}) + 1
     & = \int_{\cK} \frac{P_{Y|T_0=t}(y)^2}{P_{Y}(y)} dy
    = \int_{\cK} \frac{p_t^\lambda(y)^2}{p(y)} dy\,.
    \label{eq:chi_squared_divergence_integral_form_1}
\end{align}
Note that
\begin{align}
    p(y) = \int_{\cK} P_{Y|T_0=t}(y) P_{T_0}(t) dt
    = \int_{\cK} q(t) p_t^\lambda(y) dt\,,
    \label{eq:p_convolution_form_1}
\end{align}
so by putting this back into \eqref{eq:chi_squared_divergence_integral_form_1}, and taking the expectation over $T_0$, we have
\begin{align}
    \int_\cK \chi^2(P_{Y|T_0=t} \| P_{Y}) q(t) dt + 1
     & = \int_{\cK} q(t) \int_{\cK} \frac{p_t^\lambda(y)^2}{\int_{\cK} q(s) p_s^\lambda(y) ds} dy dt   \\
     & = \int_{\cK} \int_{\cK}q(t) \frac{p_t^\lambda(y)^2}{\int_{\cK} q(s) p_s^\lambda(y) ds} dt dy    \\
     & = \int_{\cK} \frac{\int_{\cK}q(t) p_t^\lambda(y)^2 dt}{\int_{\cK} q(s) p_s^\lambda(y) ds} dy\,.
    \label{eq:chi_squared_divergence_integral_form_2}
\end{align}
% Now we can go two ways from here. (the second way is to use $p_t^\lambda(y) q(t) = P_{Y,T}(y,t)$, but let's first try something else)
Next, we have
\begin{align}
    \int_{\cK} \frac{\int_{\cK}q(t) p_t^\lambda(y)^2 dt}{\int_{\cK} q(s) p_s^\lambda(y) ds} dy
    \leq
    \int_{\cK} \sup_{t\in\cK} p_t^\lambda(y)\,dy
    =
    \int_{\cK} \sup_{t\in\cK} r(y - \lambda t)\,dy
\end{align}
where we use the fact that $p_t^\lambda(y) = r(y - \lambda t)$.
By upper bounding naively the supremum by the \(L^\infty\) norm of \(r\), we get
\begin{align}
    \int_{\cK} \sup_{t\in\cK} r(y - \lambda t)\,dy
    \le \|r\|_{\infty} \cdot \mathrm{Vol}(\cK)\,.
\end{align}
So we just need to upper bound the \(L^\infty\) norm of \(r\), which is a scaled version of \(p\).
From the relation $r(z) = \frac{1}{(1-\lambda)^d} p\left(\frac{z}{1-\lambda}\right)$, we have
\begin{align}
    \|r\|_{\infty} = \frac{1}{(1-\lambda)^d} \|p\|_{\infty}\,.
\end{align}
Further, recall the convolution form of $p$ in \eqref{eq:p_convolution_form_1}, we can write
\begin{align}
    p(y)
    = \int_{\cK} q(t') r(y - \lambda t') dt'
    = \int_{\lambda \cK} \frac{q(t)}{\lambda^d} r(y - t) dt
    = (q_\lambda * r)(y),
\end{align}
where
\begin{align}
    q_\lambda(t) :=\begin{cases}
                       \frac{1}{\lambda^d} q\left(\frac{t}{\lambda}\right) & t \in \lambda \cK   \\
                       0                                                   & \text{otherwise}\,.
                   \end{cases}
\end{align}
By Young's convolution inequality, we have
\begin{align}
    \|p\|_{\infty} \le \|q_\lambda\|_{\infty} \|r\|_{1} = \|q_\lambda\|_{\infty} = \frac{1}{\lambda^d} \|q\|_{\infty}\,,
\end{align}
since \(r\) is a probability density.
Combining the last two displays, we have
\begin{align}
    \|r\|_{\infty} = \frac{1}{(1-\lambda)^d} \|p\|_{\infty} \le \frac{1}{\lambda (1-\lambda)^d} \|q\|_{\infty}\,,
\end{align}
and therefore,
\begin{align}
    I_{\chi^2}(T_0; Y) + 1
    \le \frac{\mathrm{Vol}(\cK)}{\lambda^d (1-\lambda)^d} \|q\|_{\infty}\,.
\end{align}
Furthermore, from the original objective, we have a factor of $\frac{1}{\lambda^2}$ in front of the chi-squared mutual information, so we have
\begin{align}
    \frac{1}{\lambda^2} I_{\chi^2}(T_0; Y)
    \le \frac{\mathrm{Vol}(\cK)}{\lambda^{d+2} (1-\lambda)^d} \|q\|_{\infty},
\end{align}
which by setting $\lambda = \frac{d + 2}{2d + 2}$ gives
\begin{align}
    \frac{1}{\lambda^2} I_{\chi^2}(T_0; Y)
    \le \frac{\mathrm{Vol}(\cK)}{\left(\frac{d + 2}{2d + 2}\right)^{d+2} \left(\frac{d}{2d + 2}\right)^d} \|q\|_{\infty}
    \le \frac{\mathrm{Vol}(\cK)}{\left(\frac{d}{2d + 2}\right)^{2d+2}} \|q\|_{\infty}
    \le \mathrm{Vol}(\cK) \cdot 4^{d+1} \|q\|_{\infty}\,.
\end{align}
\begin{remark}
    Couple of remarks on this bound:
    \begin{itemize}
        \item It is not a very good bound, since it has an exponential dependence on $d$.
        \item Also note that we got this far without using the log-concavity or any other structural assumption on $q$.
        \item It is affine invariant, since the right hand side does not change if we apply an affine transformation to $q$.
        \item $\|q\|_{\infty}$ is a very bad quantity to upper bound, since as the algorithm runs, the distribution $q$ will get more and more spiky, and consequently $\|q\|_{\infty}$ will get larger and larger.
    \end{itemize}
\end{remark}

\paragraph{Upper bounding $\|q\|_\infty$ for l.c. measures}
A possible way to upper bound $\|q\|_{\infty}$ is to use the fact that $q$ is log-concave then (using the isotropic constant / slicing bound) we have
\begin{align}
    \|q\|_{\infty} \le \frac{C^d}{\sqrt{\det(\mathrm{Cov}(q))}}\,,
\end{align}
where $C$ is some universal constant.
By putting this back into the previous bound, we have
\begin{align}
    I_{\chi^2}(T_0; Y) + 1
    \le
    \frac{\mathrm{Vol}(\cK)}{\lambda^d (1-\lambda)^d} \cdot \frac{C^d}{\sqrt{\det(\mathrm{Cov}(q))}}
    =
    \frac{C^d}{\lambda^d (1-\lambda)^d} \cdot \frac{\mathrm{Vol}(\cK)}{\sqrt{\det(\mathrm{Cov}(q))}}
\end{align}
A nice property of this upper bound is that it is \textbf{affine invariant},
i.e. it does not change if we apply an affine transformation to the distribution $q$.
A bad property of this upper bound is that the first term has an exponential dependence on $d$, which is not good for our purposes.


\clearpage
\subsubsection{An upper bound based on Fisher information matrix (WIP)}
Since $q$ is log-concave, then $p(y) = e^{-V(y)}$ for some convex potential function $V:\Rd \to \R$.
We have the autoregressive model
\begin{align}
    Y & = \lambda T_0 + (1-\lambda) \lambda T_1 + (1-\lambda)^2 \lambda T_2 + \cdots = \lambda T_0 + (1-\lambda) Y'\,,
\end{align}
where $Y, Y' \stackrel{i.i.d}{\sim} p$.
Same as the previous section, we define the density of the conditional distribution of $Y$ given $T_0 = t$ as
\begin{align}
    p^\lambda_t(y) = \frac{1}{(1-\lambda)^d} p\left(\frac{y - \lambda t}{1-\lambda}\right)
\end{align}

We want to evaluate the chi-squared mutual information quantity, which from \eqref{eq:chi_squared_divergence_integral_form_1} can be written as
\begin{align}
    I_{\chi^2}(T_0; Y) = \mathbb{E}_{T \sim q} \left[ \int_{\cK} p(y) \left( \frac{p^\lambda_T(y)}{p(y)} - 1 \right)^2 dy \right]\,,
    \label{eq:chi_squared_divergence_integral_form_3}
\end{align}
\paragraph{Taylor expansion of $\frac{p_t^\lambda(y)}{p(y)}$ around $\lambda=0$.}
Let
\begin{align}
    h(\lambda) = \frac{p_t^\lambda(y)}{p(y)} = \frac{1}{(1-\lambda)^d} \cdot \frac{p\left(\frac{y - \lambda t}{1-\lambda}\right)}{p(y)}
    = \frac{1}{(1-\lambda)^d} \cdot \frac{e^{-V\left(\frac{y - \lambda t}{1-\lambda}\right)}}{e^{-V(y)}}\,.
\end{align}
To compute $h'$, we use the fact that $h'(\lambda) = h(\lambda) \cdot \frac{d}{d\lambda} \log h(\lambda)$.
First, we compute $\log h(\lambda)$ as
\begin{align}
    \log(h(\lambda)) = -d \log(1-\lambda) + V(y) - V\left(\frac{y - \lambda t}{1-\lambda}\right)\,,
\end{align}
so
\begin{align}
    (\log h)'(\lambda)
     & = \frac{d}{1-\lambda} + \nabla V\left(\frac{y - \lambda t}{1-\lambda}\right)^\top \cdot \frac{t-y}{(1-\lambda)^2}\,,
\end{align}
and therefore
\begin{align}
    h'(\lambda)
     & = h(\lambda) \cdot \left(\frac{d}{1-\lambda} + \nabla V\left(\frac{y - \lambda t}{1-\lambda}\right)^\top \cdot \frac{t-y}{(1-\lambda)^2}\right)\,,
\end{align}
which for $\lambda = 0$ simplifies to $h'(0) = d - \langle \nabla V(y), y-t \rangle$.
Before computing $h''(\lambda)$, we can already write the Taylor expansion of $h$ around $\lambda = 0$ as
\begin{align}
    h(\lambda) = h(0) + \lambda h'(0) + \frac{1}{2} \lambda^2 h''(\tilde{\lambda}) = 1 + \lambda \Big( d - \langle \nabla V(y), y-t \rangle \Big) + \frac{1}{2} \lambda^2 h''(\tilde{\lambda})\,,
\end{align}
where $\tilde{\lambda} \in (0, \lambda)$.
By overloading a bit the notation, let $h_{t,y}(\lambda) := h(\lambda)$ to emphasize the dependence of $h$ on $t$ and $y$, then
putting this back into \eqref{eq:chi_squared_divergence_integral_form_3}, we have
\begin{align}
    I_{\chi^2}(T_0; Y) = \mathbb{E}_{T \sim q}
    \left[
    \int_{\cK} p(y) \Big( \lambda (d - \langle \nabla V(y), y-T \rangle) + \frac{\lambda^2}{2} h''_{T,y}(\tilde{\lambda}_{T,y}) \Big)^2 dy
    \right]\,,
\end{align}
\textbf{We drop the second order term for now.}

\noindent
By plugging this approximation back into the integral, the $\chi^2$ divergence becomes:
\begin{align}
    I_{\chi^2}(T_0; Y) \approx \lambda^2 \mathbb{E}_{T \sim q}
    \left[
        \int_{\cK} p(y) \Big( d - \langle \nabla V(y), y-T \rangle \Big)^2 dy
        \right]
\end{align}

We expand the squared term inside the integral. Assume without loss of generality that $q$ is centered ($\mathbb{E}[T] = 0$), which implies $Y$ is also centered.
We can split the term as $(d + \langle \nabla V(y), y \rangle) - \langle \nabla V(y), T \rangle$.
Integrating against $p(y)$ and taking the expectation gives
\begin{align}
    I_{\chi^2}(T_0; Y) \approx \lambda^2 \left(
    \underbrace{\int_{\cK} p(y) \big( d - \langle \nabla V(y), y \rangle \big)^2 dy}_{\text{First term}} + \underbrace{\int_{\cK} p(y) \langle \nabla V(y), \mathbb{E}[T T^\top] \nabla V(y) \rangle dy}_{\text{Second term}} \right)
\end{align}

\paragraph{Bounding the second term.}
Notice that $\mathbb{E}[TT^\top] = \Sigma_q$. Furthermore, the second term is exactly the trace of the Fisher Information Matrix $J(p)$ weighted by $\Sigma_q$ as
\begin{align}
    \int_{\cK} p(y) (\nabla V(y))^\top \Sigma_q (\nabla V(y)) dy = \text{Tr}(\Sigma_q J(p))\,.
\end{align}

By the stationary condition $Y = \lambda T_0 + (1-\lambda)X$, the covariance of $p$ is $\Sigma_p = \frac{\lambda}{2-\lambda} \Sigma_q$.
For a strictly log-concave distribution, the Fisher Information Matrix is bounded by a constant multiple of its inverse covariance matrix (not really..)
\begin{align}
    J(p) \preceq C \Sigma_p^{-1} = C \left( \frac{2-\lambda}{\lambda} \right) \Sigma_q^{-1} \approx \frac{2C}{\lambda} \Sigma_q^{-1}\,,
\end{align}
Taking the trace:
\begin{align}
    \text{Tr}(\Sigma_q J(p)) \le \text{Tr}\left( \Sigma_q \frac{2C}{\lambda} \Sigma_q^{-1} \right) = \frac{2C}{\lambda} \text{Tr}(I_d) = O\left(\frac{d}{\lambda}\right)
\end{align}

\paragraph{Bounding the first term.}
Let
\begin{align}
    g(Y) = d - \langle \nabla V(Y), Y \rangle, \qquad \text{and} \qquad p(y) = e^{-V(y)}\,,
\end{align}
and we want to evaluate $\E[g(Y)^2]$.
First, from integration by parts, we have
\begin{align}
    \E[g(Y)] = d - \int_{\cK} p(y) \nabla V(y)^\top y dy = d - \int_\cK \nabla p(y)^\top y dy
\end{align}
\begin{align}
    \E[g(Y)^2] = \E\left[
        (d - \nabla V(Y)^\top Y)^2
        \right]
    = \E\left[
        d^2 - 2d\nabla V(Y)^\top Y + Y^\top \nabla V(Y) \nabla V(Y)^\top Y
        \right]
\end{align}
Meanwhile, the term $\E_Y[g(Y)^2] = \text{Var}(\langle \nabla V(Y), Y \rangle)$ represents the variance of the trace of the Hessian-equivalent, which for standard log-concave measures scales like $O(d^2)$ (this also needs to be verified).

Finally, substitute these bounds back into the target quantity:
\begin{align}
    \frac{1}{\lambda^2} I_{\chi^2}(T_0; Y) \approx \frac{1}{\lambda^2} \left( \lambda^2 O(d^2) + \lambda^2 O\left(\frac{d}{\lambda}\right) \right) = O(d^2) + O\left(\frac{d}{\lambda}\right)
\end{align}
To evaluate $\inf_{\lambda \in (0,1)} \frac{1}{\lambda^2} I_{\chi^2}(T_0; Y)$, we are free to choose a scaling for $\lambda$. If we set $\lambda = \frac{c}{d}$ (the exact same scaling that optimizes the exact Gaussian case), we obtain:$$\frac{1}{\lambda^2} I_{\chi^2}(T_0; Y) \le O(d^2) + O\left( d \cdot \frac{d}{c} \right) = O(d^2)$$Thus, the supremum over any compact log-concave distribution $q$ yields a polynomial upper bound of $O(d^2)$, matching the dimensional scaling of the Gaussian case.

\newpage
Let $q$ be a log-concave distribution on a convex body $\cK\subseteq \Rd$, and let $T_k$ be i.i.d. samples from $q$.
Does there exist a $\poly(d)$ upper bound on the following chi-squared mutual information quantity?
\begin{align}
    \sup_{q} \inf_{\lambda\in(0,1)} \frac{1}{\lambda^2} I_{\chi^2}\left(T_0; \lambda \sum_{k=0}^\infty (1-\lambda)^k T_k\right)
\end{align}

\begin{example}[The Gaussian case]
    Let
    \[
        Y \triangleq \lambda \sum_{k=0}^\infty (1-\lambda)^k T_k
    \]
    and let $p$ be the law of $Y$.

    If $q=\cN(\mu,\Sigma_q)$ on $\R^d$, then the affine recursion
    $X_{n+1}=(1-\lambda)X_n+\lambda T_{n+1}$ has the unique stationary law
    \[
        p=\cN(\mu,\Sigma_p),
        \qquad
        \Sigma_p=\frac{\lambda}{2-\lambda}\Sigma_q.
    \]
    Moreover, for fixed $t\in\R^d$, the conditional law $\mu_t$ of
    \[
        Y=(1-\lambda)X+\lambda t,\qquad X\sim p
    \]
    is Gaussian with mean $(1-\lambda)\mu+\lambda t$ and covariance $(1-\lambda)^2\Sigma_p$.
    The pair $(T,Y)$ is jointly Gaussian, and in whitened coordinates it splits into $d$
    independent correlated Gaussian pairs with
    \[
        \rho^2=\lambda(2-\lambda)
        \qquad\Longrightarrow\qquad
        1-\rho^2=(1-\lambda)^2.
    \]
    Using the chi-squared mutual information formula for jointly Gaussian variables,
    \[
        I_{\chi^2}(T;Y)=\chi^2(P_{T,Y}\|P_TP_Y)
        = \prod_{i=1}^d\frac{1}{1-\rho^2}-1
        = (1-\lambda)^{-2d}-1.
    \]
    Therefore,
    \[
        \frac{1}{\lambda^2}I_{\chi^2}(T;Y)
        = \frac{1}{\lambda^2}\bigl((1-\lambda)^{-2d}-1\bigr).
    \]
    In particular, taking $\lambda=\frac{c}{d}$ gives
    \[
        (1-\lambda)^{-2d}-1=\left(1-\frac{c}{d}\right)^{-2d}-1 \approx e^{2c}-1
        \,\Rightarrow\,
        \sup_{q} \inf_\lambda I_{\chi^2}(T;Y) \le \frac{e^{2c}-1}{c^2}\,d^2.
    \]
\end{example}


\clearpage
\bibliographystyle{plainnat}
\bibliography{../all}
\clearpage

\end{document}